\chapter*{Balanç de la part teòrica}
\phantomsection
\addcontentsline{toc}{chapter}{Balanç de la part teòrica}
\label{cap:balanc}

Els capítols teòrics estableixen la infraestructura de la teoria estructural de conjunts, des del llenguatge d'aplicacions fins a la comparació general de mides i el llenguatge de les famílies indexades.

\begin{itemize}
  \item El capítol~\ref{cap:aplicacions} introdueix el llenguatge d'aplicacions, dominis, codominis i composició; defineix els conjunts com a identitats; i fixa els conjunts terminal i inicial, els elements generalitzats, els elements globals i la separació global.
  \item El capítol~\ref{cap:parts} desenvolupa monomorfismes, epimorfismes, parts i pertinença generalitzada. També introdueix sumes, productes, equalitzadors, pullbacks, imatges inverses i interseccions.
  \item El capítol~\ref{cap:valors} reforça les sumes amb la disjunció i l'estabilitat, distingeix els dos valors de veritat, construeix el complement i les operacions booleanes amb parts, i obté tots els límits i colímits finits.
  \item El capítol~\ref{cap:relacions} desenvolupa les relacions sobre elements generalitzats, les relacions d'equivalència en forma estructural, el parell nucli, els quocients com a coequalitzadors, els epimorfismes regulars i les particions, units en la correspondència de l'axioma d'efectivitat; i, amb l'axioma d'elecció, estableix que tot epimorfisme és escindit i regular.
  \item El capítol~\ref{cap:imatges} construeix la factorització de tota aplicació com un epimorfisme regular seguit d'un monomorfisme, defineix la imatge directa i la inversa amb la seva adjunció, demostra l'estabilitat dels epimorfismes regulars per canvi de base ---amb la qual la teoria disposa de la infraestructura d'una categoria regular--- i en treu la composició de relacions i les relacions funcionals.
  \item El capítol~\ref{cap:exponencials} postula l'existència d'exponencials, estableix la propietat universal de transposició i la unicitat de $B^A$, en calcula els casos amb $\mathbf0$, $\mathbf1$ i $\mathbf2$ ---aquest últim definint el conjunt de parts $P(A):=\mathbf2^A$---, i realitza internament la composició i les famílies d'aplicacions.
  \item El capítol~\ref{cap:parts-interna} desenvolupa el conjunt de parts com a objecte: la relació de pertinença interna, la identificació de les aplicacions cap a $P(A)$ amb les parts parametritzades, les operacions booleanes internes amb el seu càlcul puntual, la inclusió interna com a ordre parcial, i la imatge inversa, la imatge directa i el quantificador universal com a cadena d'adjuncions entre conjunts de parts.
  \item El capítol~\ref{cap:naturals} postula l'objecte dels nombres naturals mitjançant la propietat universal de recursió, en deriva el teorema de recursió, el principi d'inducció i les equacions de Peano, i construeix l'aritmètica ---suma, producte i exponenciació--- amb la recursió amb paràmetres via els exponencials.
  \item El capítol~\ref{cap:finits} converteix la finitud en una noció interna: els ordinals finits $[n]$ per recursió, el conjunt finit com el que és isomorf a algun $[n]$, la unicitat del cardinal, el principi dels calaixos i l'equivalència ---pròpia del cas finit--- entre injectivitat, exhaustivitat i bijectivitat, i el recompte de parts $2^n$ i d'aplicacions $m^n$.
  \item El capítol~\ref{cap:infinitud} estén la comparació de cardinals més enllà del cas finit: els conjunts infinits i els infinits de Dedekind i la seva relació ---amb l'elecció numerable com a bisagra---, els conjunts enumerables i numerables i la seva estabilitat sense elecció, els subobjectes de $\N$ com a sempre enumerables gràcies al mínim, les llistes finites amb el seu principi de recursió, i les unions indexades numerablement, on l'elecció numerable torna a ser essencial.
  \item El capítol~\ref{cap:cantor} obre la comparació general de mides: la relació de dominació per monomorfismes, el teorema de Schröder--Bernstein pel menor punt fix d'un operador creixent sobre $P(A)$ ---sense elecció---, l'argument diagonal de Cantor i la desigualtat estricta $A\prec P(A)$, amb la conseqüència que no hi ha conjunt universal ni mida màxima.
  \item El capítol~\ref{cap:families} fixa el llenguatge de les famílies indexades: una família és una aplicació $p\colon E\to I$ amb les fibres com a producte fibrat, les seccions com a funcions d'elecció, la suma i el producte d'una família, les famílies de parts d'un fix amb les seves unions i interseccions indexades, i la distinció central entre una família interna existent i una regla externa encara no representada, amb la substitució estructural SR enunciada com a esquema.
  \item El capítol~\ref{cap:bons-ordres} desenvolupa els ordres, els bons ordres i els segments inicials: el bon ordre com \enquote{tota part no inicial té mínim}, la rigidesa ---entre dos bons ordres hi ha com a molt un isomorfisme d'ordre--- i la comparació per segments, amb la tricotomia establerta en la seva part d'exclusivitat.
  \item El capítol~\ref{cap:recursio} estableix la inducció i la recursió transfinites amb codomini fix ---la reunió de solucions parcials dins un exponencial fix, sense elecció--- i en treu la comparabilitat total dels bons ordres.
  \item El capítol~\ref{cap:zorn} reuneix les tres formes de maximalitat: el teorema de Hartogs (un bon ordre que no s'injecta en $A$), el teorema del bon ordre i el lema de Zorn, amb totes les recursions acotades per $\Hart$, i l'equivalència entre elecció, bon ordre i Zorn.
  \item El capítol~\ref{cap:ordinals} defineix els tipus ordinals com a bons ordres llevat d'isomorfisme, amb zero, successor i límit, els suprems acotats, i el teorema de Burali--Forti com a paral·lel ordinal de Cantor.
  \item El capítol~\ref{cap:aritmetica-ordinal} construeix l'aritmètica ordinal ---suma, producte i exponenciació sobre representants---, no commutativa ($1+\omega=\omega\neq\omega+1$), amb la distinció entre ordre i mida com a fil conductor.
  \item El capítol~\ref{cap:cardinals} desenvolupa els cardinals i la seva aritmètica, relacional i sense representants, amb Cantor com $\abs{A}<2^{\abs{A}}$, la comparabilitat total sota elecció, els cardinals inicials i el teorema de König.
  \item El capítol~\ref{cap:alephs}, opcional i de Nivell~II, recorre la jerarquia dels alephs i dels beths i la cofinalitat fent servir l'esquema de reunió de famílies SR, i formula amb precisió la hipòtesi del continu com una qüestió independent ---tot i que el continu $\mathfrak c=\abs{P(\N)}$ i la mateixa equació d'HC $2^{\aleph_0}=\aleph_1$ són ja de Nivell~I, perquè $P(\N)$ i $\aleph_1=(\aleph_0)^+$ existeixen sense reunir cap família no acotada; només la successió sencera dels alephs demana SR.
  \item El capítol~\ref{cap:models} tanca el volum amb un balanç metateòric: recull tots els axiomes, caracteritza l'estructura com un model d'ETCS, distingeix la teoria de Nivell~I d'ETCS reforçat amb SR, compara l'ordinal material de von Neumann amb el tipus estructural, i en delimita els límits de mida, consistència i força expressiva.
\end{itemize}

En el metallenguatge, les classes d'equivalència de parts d'un conjunt formen una àlgebra booleana, i les relacions d'equivalència es corresponen amb les particions; internament, aquesta lògica de subconjunts es realitza com l'àlgebra d'un únic objecte $P(A)$, amb la pertinença, les operacions booleanes, l'ordre d'inclusió i els quantificadors ja construïts. L'objecte dels nombres naturals proporciona la recursió, la inducció i l'aritmètica, sobre les quals es construeixen la finitud interna, el recompte i la cardinalitat numerable, amb el paper de l'elecció numerable delimitat amb precisió. La segona part porta la construcció fins a la teoria clàssica de l'infinit: Schröder--Bernstein i Cantor obren la comparació general de mides; els bons ordres i la recursió transfinita, amb Hartogs i el teorema del bon ordre, donen la teoria de l'ordre; els tipus ordinals i els cardinals, amb les seves aritmètiques, separen la forma d'ordre de la mida; i Burali--Forti fa de paral·lel ordinal de Cantor. Tot això es fa amb els axiomes de Nivell~I ---o amb l'elecció---, mantenint les construccions acotades; només la reunió de les col·leccions transfinites no acotades ---la successió sencera dels alephs i dels beths--- demana l'esquema de substitució estructural SR, que porta la teoria a la força de ZFC. (El continu mateix, $\mathfrak c=\abs{P(\N)}$, i la hipòtesi del continu són de Nivell~I: no depenen de SR, sinó només del conjunt de parts i del cardinal successor per Hartogs.) El balanç metateòric final situa tot l'edifici, el caracteritza com un model d'ETCS, en distingeix la versió de Nivell~I de la reforçada amb SR, i el compara amb la teoria clàssica de von Neumann.
